% unit: noether.p10.title.001
\begin{center}
{\Large\bfseries 10. Les équations fonctionnelles de l'application\\ isomorphe}\par
\vspace{0.75em}
Par Emmy Noether, à Göttingen.\par
\vspace{1em}
\emph{Math. Ann.} 77 (1916), pp. 536--545
\end{center}

\vspace{2em}

% unit: noether.p10.par.001
D'après Dedekind,\footnote{Cf. par exemple Dirichlet-Dedekind, \emph{Leçons sur la théorie des nombres} (4e éd.), Supplément XI, \S\ 161. La désignation \og{}application isomorphe\fg{} ou \og{}isomorphisme\fg{}, calquée sur la théorie des groupes, ne se rencontre que dans la littérature postérieure; Dedekind parle des \og{}permutations d'un corps\fg{}.} on entend par \emph{application isomorphe du corps} $\Afield$ \emph{sur un système} $\Bfield$ --- qui se révèle après coup être lui aussi un corps --- \emph{une correspondance univoque telle qu'à chaque élément de} $\Afield$ \emph{corresponde un et un seul élément de} $\Bfield$; \emph{et qu'à la somme, à la différence, au produit et au quotient de deux éléments quelconques de} $\Afield$ \emph{soient de nouveau associés la somme, la différence, le produit et le quotient des éléments de} $\Bfield$ \emph{qui leur correspondent univoquement}.

% unit: noether.p10.par.002
Une telle application est réalisée par une fonction $f(z)$, univoque d'après ce qui précède. Si alors $z$ parcourt tous les éléments $x,y,\ldots$ du corps $\Afield$, $f(z)$ est donc caractérisée par les équations fonctionnelles
% unit: noether.p10.eq.001
\[
\begin{array}{ll}
(1)\quad f(x+y)=f(x)+f(y); & (2)\quad f(x-y)=f(x)-f(y);\\[0.45em]
(3)\quad f(x\cdot y)=f(x)\cdot f(y); & (4)\quad f\!\left(\dfrac{x}{y}\right)=\dfrac{f(x)}{f(y)}.
\end{array}
\]

% unit: noether.p10.par.003
\emph{Dans ce qui suit seront indiquées les solutions univoques les plus générales} $f(z)$ \emph{de ces équations fonctionnelles, lorsque} $x,y,\ldots$ \emph{parcourent tous les nombres réels et complexes; c'est-à-dire la fonction la plus générale} $f(z)$ \emph{qui réalise une application isomorphe du corps de tous les nombres complexes}.\footnote{Cette question m'a été posée occasionnellement par M. Landau. Comme je l'ai remarqué après coup, une partie des solutions se trouve déjà chez H. Lebesgue, \emph{Sur les transformations ponctuelles, transformant les plans en plans}, Atti Acad. Torino, 1906/07; et aussi implicitement chez A. Ostrowski, \emph{Über einige Fragen der allgemeinen Körpertheorie}, Crelle's Journal 143 (1913), \S\ 2.} La solution pour des corps arbitraires s'obtient de façon tout à fait analogue.

% unit: noether.p10.par.004
G. Hamel\footnote{G. Hamel, \emph{Eine Basis aller Zahlen und die unstetigen Lösungen der Funktionalgleichung:} $f(x+y)=f(x)+f(y)$, Math. Ann. 60, p. 459 (1905).} a construit les solutions les plus générales de l'équation fonctionnelle (1) au moyen d'une base \emph{linéaire} de tous les nombres. La construction donnée ici pour les solutions de (1) à (4) --- donc pour la fonction d'application relativement à laquelle les opérations rationnelles et algébriques sont invariantes --- utilise la base \emph{rationnelle} et \emph{algébrique} de tous les nombres. Après que l'application isomorphe aura été discutée de plus près en 1, des conditions nécessaires pour $f(z)$ seront données en 2; en 3 elles seront reconnues comme suffisantes et conduiront à la construction effective.\footnote{Cf. les considérations parallèles dans mon travail: \emph{Die allgemeinsten Bereiche aus ganzen transzendenten Zahlen}. Math. Ann. 77, p. 103 (1915), en particulier \S\S\ 8 et 10.} Les solutions de (1) à (4), à l'exception connue de $f(z)=z$ ou $\overline z$, sont discontinues; en 4 on montre en outre qu'elles sont \og{}extrêmement discontinues\fg{}, fait que G. Hamel a démontré pour les solutions réelles de (1), mais qui n'a pas nécessairement lieu dans le domaine complexe.\footnote{L'expression de Hamel \og{}totalement discontinu\fg{} est employée dans le reste de la littérature en un autre sens.}

% unit: noether.p10.sec.001
\medskip
\noindent\textbf{1.}\quad Nous tirons d'abord quelques conséquences des équations fonctionnelles (1) à (4), directement pour des corps généraux $\Afield$. D'après (1), l'univocité entraîne $f(0)=0$. Si $f(y)$ ne s'annule pas, alors le $y$ initial ne peut pas non plus s'annuler, et les équations fonctionnelles (1) à (4) montrent ainsi que \emph{le système d'application} $\Bfield$ \emph{de toutes les valeurs} $f(z)$ \emph{forme de nouveau un corps}. Inversement, si l'on prescrit la condition initiale $f(0)=0$, alors (2) entraîne l'univocité de la fonction $f(z)$: \emph{l'exigence d'univocité et la condition initiale} $f(0)=0$ \emph{sont donc équivalentes}. De (4) il résulte que $f(z)$ ne peut pas s'annuler identiquement; de (3), en outre, que $f(z)$ ne s'annule que pour $z=0$. Car de $f(y)=0$ et $y\ne0$ il résulterait, puisque $z/y$ appartient aussi à $\Afield$, que
% unit: noether.p10.eq.002
\[
f(z)=f\!\left(\frac{z}{y}\cdot y\right)=f\!\left(\frac{z}{y}\right)\cdot f(y)=0,
\]
ce qui contredirait (4) en donnant l'annulation identique de $f(z)$. L'égalité $f(x)=f(y)$ entraîne donc $x=y$; c'est-à-dire qu'à chaque valeur de $f(z)$ correspond aussi une et une seule valeur de $z$: \emph{$f(z)$ est une fonction réciproquement univoque de} $z$. Comme le montrent les équations fonctionnelles, l'application définie par la fonction inverse est à son tour isomorphe. Puisque toute opération rationnelle est composée d'un nombre fini d'additions, de multiplications et de leurs inverses, on peut donc dire aussi: \emph{par l'application isomorphe est établie une relation biunivoque entre les éléments de} $\Afield$ \emph{et ceux de} $\Bfield$, \emph{de telle sorte que toutes les relations rationnelles existant entre un nombre fini quelconque d'éléments de} $\Afield$,
% unit: noether.p10.eq.003
\[
F(x,y,\ldots,t)=0,
\]
\emph{sont conservées dans} $\Bfield$, \emph{et réciproquement}.

% unit: noether.p10.par.005
Il faut encore remarquer que, pour un corps, $f(z)$ est déjà caractérisée par les équations fonctionnelles (1) et (3), par l'exigence que $f(z)$ ne s'annule pas identiquement, et par la condition initiale $f(0)=0$. En effet, (2) résulte de (1) en remplaçant $x$ par l'élément $(x-y)$ appartenant à $\Afield$; puis (2) et $f(0)=0$ donnent l'univocité de $f(z)$. Comme on l'a montré plus haut, il résulte en outre de (3), en remplaçant $x$ par l'élément $x/y$ appartenant à $\Afield$, que $f(y)$ ne peut pas s'annuler pour $y\ne0$, donc la validité de (4) et en même temps la biunivocité. Si, au lieu d'un corps, on prend pour base un domaine d'intégrité $\Sdomain$, $x/y$ n'appartient pas nécessairement à $\Sdomain$, et l'on ne peut pas conclure de (3) à la biunivocité. Ici suffit cependant la formulation la plus connue: une application isomorphe est une correspondance \emph{biunivoque} entre les éléments de $\Sdomain$ et de $\Sdomain'$ telle qu'à la somme et au produit correspondent toujours la somme et le produit.\footnote{Cf. pour le n\textsuperscript{o} 1: Dedekind, loc. cit., \S\ 161.}

% unit: noether.p10.sec.002
\medskip
\noindent\textbf{2.}\quad Désormais, pour simplifier, on prendra pour base, à la place de $\Afield$, le corps $\Cfield$ de tous les nombres réels et complexes. Il s'agit d'abord de discuter les systèmes de valeurs que prend $f(z)$. De (3) ou de (4), on déduit que $f(1)=1$, et dès lors, en vertu des équations fonctionnelles, pour tout \emph{nombre rationnel} $\alpha$, $f(\alpha)=\alpha$; \emph{ainsi tout nombre rationnel correspond à lui-même dans l'application}. Comme un \emph{nombre algébrique} $\beta$ satisfait à une équation irréductible à coefficients numériques rationnels $F(\beta,\alpha)=0$, le n\textsuperscript{o} 1 et la remarque précédente donnent, pour $f(\beta)$, l'équation $F(f(\beta),\alpha)=0$; \emph{tout nombre algébrique passe donc en lui-même ou en l'un de ses conjugués}, et, par biunivocité, à la totalité des nombres algébriques correspond de nouveau le corps de tous les nombres algébriques. Afin de reconnaître le comportement des nombres \emph{transcendants}, et en même temps de séparer les valeurs conjuguées, nous prenons pour base une \emph{base rationnelle de tous les nombres},\footnote{Cf. mon travail cité sur les \og{}nombres transcendants entiers\fg{}, \S\ 6.} c'est-à-dire un système $\Theta$ de nombres $\vartheta$ qui permet d'exprimer tous les nombres rationnellement --- avec des coefficients numériques rationnels\footnote{Par \og{}fonction rationnelle\fg{}, j'entendrai toujours une fonction dont les coefficients sont eux aussi des nombres rationnels.} --- tandis que, dans au moins un bon ordre, aucun nombre de base n'est exprimable rationnellement au moyen d'un nombre fini de précédents. L'existence de cette base résulte du théorème du bon ordre exactement comme celle des bases linéaires et algébriques.

% unit: noether.p10.par.006
Soumettons $\Cfield$ à un bon ordre $\Omega$. Alors tout nombre est ou bien exprimable rationnellement au moyen d'un nombre fini de nombres précédents, ou bien rationnellement indépendant des précédents. Ces derniers nombres forment la base $\Theta$, et le raisonnement par induction montre qu'effectivement tout nombre est exprimable rationnellement au moyen d'un nombre fini de $\vartheta$.

% unit: noether.p10.par.007
À chaque nombre de base $\vartheta$ est défini le \og{}corps de segment $\Kfield(\vartheta)$\fg{}: il est constitué de toutes les combinaisons rationnelles des nombres de base qui précèdent $\vartheta$ dans le bon ordre. Comme corps de segment du premier nombre de base, il faut considérer le corps $\Rfield$ de tous les nombres rationnels. Le nombre $\vartheta$ peut présenter deux sortes de comportement par rapport à $\Kfield(\vartheta)$: il peut être algébriquement dépendant de $\Kfield(\vartheta)$, ou algébriquement indépendant de lui. Puisque, d'après le n\textsuperscript{o} 1, toutes les relations valables dans $\Cfield$ sont aussi valables dans le corps image, et réciproquement, la totalité des valeurs $f(\vartheta)$ correspondant aux $\vartheta$ forme une base du domaine image, bien ordonnée au moyen du bon ordre de $\Theta$ lui-même. En particulier, au corps de segment $\Kfield(\vartheta)$ correspond le corps de segment $\Kfield(f(\vartheta))$; et selon que $\vartheta$ est algébriquement dépendant ou indépendant de $\Kfield(\vartheta)$, il en va de même de $f(\vartheta)$ relativement à $\Kfield(f(\vartheta))$. Dans le cas de dépendance, les équations irréductibles pour $\vartheta$ et pour $f(\vartheta)$ se correspondent aussi. La totalité des valeurs $\vartheta$ qui sont à chaque fois algébriquement indépendantes de $\Kfield(\vartheta)$ forment --- comme le montre encore l'induction --- une \emph{base algébrique de tous les nombres},\footnote{Cf. E. Zermelo, \emph{Über ganze transzendente Zahlen}, \S\ 1, Math. Ann. 75 (1914).} c'est-à-dire un système $H$ de nombres $\eta$ qui permet d'exprimer algébriquement tous les nombres, tandis qu'aucune relation algébrique n'existe entre les $\eta$ eux-mêmes. À cette base algébrique $H$ correspond dans le domaine image de nouveau une base algébrique $Z$, qui, par biunivocité, a la même puissance que $H$, à savoir la puissance du continu, puisque l'on sait qu'une extension algébrique n'augmente pas la puissance.\footnote{Cf. E. Steinitz, \emph{Algebraische Theorie der Körper}, Crelle's Journal 137 (1910), \S\ 24.} Mais de la connaissance des valeurs $f(\vartheta)$ on obtient immédiatement $f(z)$ pour tous les systèmes de valeurs, puisque l'égalité $z=\psi(\vartheta)$ entraîne toujours $f(z)=\psi(f(\vartheta))$.

% unit: noether.p10.sec.003
\medskip
\noindent\textbf{3.}\quad Nous montrons maintenant que les conditions nécessaires trouvées ci-dessus fournissent effectivement aussi la construction de $f(z)$. La base $Z$ qui correspond biunivoquement à la base algébrique $H$ est construite et associée à $H$ de la manière suivante. Posons un second bon ordre $\Omega'$, qui fournit une base algébrique $Z'$ de tous les nombres. Que $Z$, ayant pour éléments les $\xi$, soit un sous-ensemble de $Z'$ de même puissance que $Z'$, et donc que $H$. À chaque élément $\eta_i$ de $H$, associons un et un seul élément $\xi_i$, portant la même numérotation, par la règle suivante: $\xi_i$ est le \emph{premier}, dans le bon ordre $\Omega'$, de tous les éléments de $Z$ qui n'ont été associés à aucun élément de $H$ précédant $\eta_i$.

% unit: noether.p10.par.008
Définissons maintenant $f(z)$ ainsi:
\begin{enumerate}
\item[(a)] \emph{pour tous les nombres rationnels} $\alpha$, \emph{par} $f(0)=0$, $f(\alpha)=\alpha$;
\item[(b)] \emph{pour toutes les grandeurs de la base algébrique} $H$, \emph{par} $f(\eta_i)=\xi_i$;
\item[(c)] \emph{pour toutes les autres grandeurs} $\vartheta$ \emph{de la base rationnelle} $\Theta$ --- \emph{celles qui sont algébriquement dépendantes du corps de segment} $\Kfield(\vartheta)$, \emph{et satisfont donc à une équation} $F(\vartheta;\vartheta_{i_1},\ldots,\vartheta_{i_\tau})=0$ \emph{irréductible dans} $\Kfield(\vartheta)$ --- \emph{comme racine de l'équation} $F(f(\vartheta);f(\vartheta_{i_1}),\ldots,f(\vartheta_{i_\tau}))=0$;
\item[(d)] \emph{pour toutes les autres valeurs} $z$, \emph{qui, par la définition de} $\Theta$, \emph{sont représentables sous la forme} $z=\psi(\vartheta_{k_1},\ldots,\vartheta_{k_\sigma})$, \emph{par} $f(z)=\psi(f(\vartheta_{k_1}),\ldots,f(\vartheta_{k_\sigma}))$.\footnote{Dans (d), (a) est contenu comme cas particulier.}
\end{enumerate}

% unit: noether.p10.par.009
Pour montrer que le $f(z)$ ainsi défini satisfait effectivement aux équations fonctionnelles (1) à (4), il suffit, d'après le n\textsuperscript{o} 1, de le prouver pour (1) et (3), puisque $\Cfield$ est un corps et que $f(z)$, à cause de (a) et (b), ne peut pas s'annuler identiquement et satisfait en outre à la condition initiale $f(0)=0$. Au moyen de la base rationnelle $\Theta$, exprimons $x,y,x+y$ sous la forme
% unit: noether.p10.eq.004
\[
x=\psi_1(\vartheta_1\cdots\vartheta_t),\quad
 y=\psi_2(\vartheta_1\cdots\vartheta_t),\quad
 (x+y)=\psi_3(\vartheta_1\cdots\vartheta_t).
\]
Ainsi la preuve que $f(z)$ satisfait à l'équation fonctionnelle (1),
% unit: noether.p10.eq.005
\[
f(x)+f(y)-f(x+y)=0,
\]
revient, d'après (d), à montrer que l'égalité
% unit: noether.p10.eq.006
\[
\psi_1(\vartheta)+\psi_2(\vartheta)-\psi_3(\vartheta)
 =\frac{H_1(\vartheta)}{H_2(\vartheta)}=0
\]
entraîne toujours
% unit: noether.p10.eq.007
\[
\psi_1(f(\vartheta))+\psi_2(f(\vartheta))-\psi_3(f(\vartheta))
 =\frac{H_1(f(\vartheta))}{H_2(f(\vartheta))}=0,
\]
où $H_1,H_2$ désignent des fonctions rationnelles entières de leurs arguments. L'équation fonctionnelle (3) conduit exactement à la même forme de condition; ainsi la satisfaction de toutes les équations fonctionnelles est identique à la condition suivante --- nécessaire aussi d'après les n\textsuperscript{os} 1 et 2: pour \emph{toute} fonction rationnelle entière $H(\vartheta)$, l'égalité $H(\vartheta)=0$ entraîne toujours $H(f(\vartheta))=0$, et $H(\vartheta)\ne0$ entraîne $H(f(\vartheta))\ne0$, ce dernier point à cause de l'apparition du dénominateur.

% unit: noether.p10.par.010
Que tel soit effectivement le cas se montre par induction. Supposons que, dans $H(\vartheta_1\cdots\vartheta_t)$, la grandeur de base $\vartheta_t$ soit, dans le bon ordre, la dernière parmi $\vartheta_1\cdots\vartheta_t$.\footnote{Les expressions $H(\vartheta)$ qui sont de degré zéro en les $\vartheta$ passent en elles-mêmes sous l'application et n'ont donc pas besoin d'être examinées davantage.} Alors $H(\vartheta_1\cdots\vartheta_t)=0$ peut être considéré comme une relation à laquelle satisfait $\vartheta_t$ relativement au corps de segment $\Kfield(\vartheta_t)$, et de même $H(f(\vartheta_1)\cdots f(\vartheta_t))=0$ comme une relation à laquelle satisfait $f(\vartheta_t)$ relativement à $\Kfield(f(\vartheta_t))$. Si toutes les relations $H(\vartheta)=0$ n'étaient pas également remplies par $f(\vartheta)$, et réciproquement, il devrait exister une première grandeur de base, disons $\vartheta_0$, pour laquelle au moins une relation valable relativement à $\Kfield(\vartheta_0)$ ne serait pas conservée dans le corps image, ou inversement, alors que les corps de segment précédents $\Kfield(\vartheta_0)$ et $\Kfield(f(\vartheta_0))$ seraient isomorphes. En particulier, d'après (a), le corps de segment de la première grandeur de base de $\Theta$, le corps $\Rfield$ de tous les nombres rationnels, est isomorphe à son corps image. Supposons maintenant que $H(\vartheta_0)$ ait la forme
% unit: noether.p10.eq.008
\[
H(\vartheta_0)=a_0(\vartheta)\cdot\vartheta_0^\chi
+a_1(\vartheta)\cdot\vartheta_0^{\chi-1}+\cdots+a_\chi(\vartheta),
\]
où les $a_i(\vartheta)$ sont des grandeurs de $\Kfield(\vartheta_0)$. Il y a alors deux possibilités.

% unit: noether.p10.par.011
1) $\vartheta_0$ est algébriquement indépendant de $\Kfield(\vartheta_0)$. Alors $H(\vartheta_0)=0$ est satisfait identiquement en $\vartheta_0$,\footnote{Comme $x,y,\ldots$ peuvent être exprimés par les $\vartheta$, cette forme de relation peut très bien se produire.} de sorte que $a_i(\vartheta)=0$ pour tous les coefficients, et donc, par l'isomorphisme de $\Kfield(\vartheta_0)$ et $\Kfield(f(\vartheta_0))$, aussi $a_i(f(\vartheta))=0$. L'égalité $H(\vartheta_0)=0$ entraîne donc toujours $H(f(\vartheta_0))=0$. Réciproquement, supposons $H(f(\vartheta_0))=0$. Puisque $\vartheta_0$, étant algébriquement indépendant de $\Kfield(\vartheta_0)$, appartient à la base algébrique $H$, $f(\vartheta_0)$ appartient par (b) à la base algébrique $Z$ et est algébriquement indépendant de $\Kfield(f(\vartheta_0))$. Ainsi tous les $a_i(f(\vartheta))$ s'annulent, et par l'isomorphisme tous les $a_i(\vartheta)$ s'annulent. L'égalité $H(f(\vartheta_0))=0$ entraîne $H(\vartheta_0)=0$; par contraposition, $H(\vartheta_0)\ne0$ entraîne aussi $H(f(\vartheta_0))\ne0$.

% unit: noether.p10.par.012
2) $\vartheta_0$ est algébriquement dépendant de $\Kfield(\vartheta_0)$. Alors $H(\vartheta_0)$ est divisible par l'équation $F(\vartheta_0)$ irréductible relativement à $\Kfield(\vartheta_0)$; c'est-à-dire que, identiquement en $t$,
% unit: noether.p10.eq.009
\[
H(t;a_i(\vartheta))=F(t;\vartheta)\cdot G(t;\vartheta_i).
\]
Comme tous les coefficients qui interviennent appartiennent à $\Kfield(\vartheta_0)$, la relation correspondante est remplie dans le corps isomorphe $\Kfield(f(\vartheta_0))$:
% unit: noether.p10.eq.010
\[
H(t;a_i(f(\vartheta)))=F(t;f(\vartheta))\cdot G(t;f(\vartheta)).
\]
Mais, d'après (c), $f(\vartheta_0)$ a été choisi comme racine de $F(t;f(\vartheta))$. L'égalité $H(\vartheta_0)=0$ entraîne donc $H(f(\vartheta_0))=0$.

% unit: noether.p10.par.013
Réciproquement, si $H(f(\vartheta_0))=0$, alors la divisibilité de $H(t;a_i(f(\vartheta)))$ par la fonction $F(t;f(\vartheta))$, irréductible relativement à $\Kfield(f(\vartheta_0))$, entraîne, pour le corps isomorphe $\Kfield(\vartheta_0)$, la divisibilité de $H(t;a_i(\vartheta))$ par $F(t;\vartheta)$. Et puisque $F(t;f(\vartheta))$ provenait de l'équation irréductible de $\vartheta_0$ par la relation isomorphe, et que cette relation est biunivoque, $F(t;\vartheta)$ représente nécessairement la fonction irréductible dont $\vartheta_0$ est racine. Ainsi, $H(f(\vartheta_0))=0$ implique $H(\vartheta_0)=0$; de même, $H(\vartheta_0)\ne0$ implique $H(f(\vartheta_0))\ne0$.

% unit: noether.p10.par.014
De corps de segment isomorphes $\Kfield(\vartheta_0)$ et $\Kfield(f(\vartheta_0))$ naissent donc aussi, par adjonction de $\vartheta_0$ et de $f(\vartheta_0)$ respectivement, des corps isomorphes; l'hypothèse que cela échoue pour un certain $\vartheta_0$ conduit à une contradiction. Il est ainsi démontré que \emph{la fonction} $f(z)$ \emph{définie par (a) à (d) représente effectivement une solution des équations fonctionnelles (1) à (4), et, d'après 2, la plus générale}.

% unit: noether.p10.par.015
Toutes les considérations utilisées pour construire $f(z)$ n'ont employé que le fait que la totalité $\Cfield$ de tous les nombres forme un corps; elles valent donc pour tout corps $\Afield$ défini abstraitement. Il faut noter ici que le corps $\Rfield$ des nombres rationnels est remplacé par le \og{}corps premier\fg{} de $\Afield$, c'est-à-dire par le corps dérivé de l'élément unité de $\Afield$ par les opérations d'addition et de multiplication et leurs inverses. Si donc, dans (a) à (d), on remplace les nombres rationnels par les éléments du corps premier, alors \emph{la fonction} $f(z)$ \emph{ainsi obtenue réalise l'application isomorphe la plus générale d'un corps arbitraire défini abstraitement}.

% unit: noether.p10.sec.004
\medskip
\noindent\textbf{4.}\quad Il reste à montrer que les fonctions $f(z)$ définies pour le corps $\Cfield$ de tous les nombres réels et complexes --- qui, à l'exception de $f(z)=z$ ou $\overline z$, sont connues pour être discontinues --- sont \emph{extrêmement discontinues}; c'est-à-dire que, dans tout voisinage de n'importe quel nombre réel ou complexe $z_0$, il existe des valeurs $z$ pour lesquelles $f(z)$ s'approche arbitrairement de n'importe quelle valeur prescrite $Z_0$. Cette extrême discontinuité appartient, comme G. Hamel l'a démontré loc. cit., aux solutions \emph{réelles} $\varphi(x)$, définies pour tous les nombres \emph{réels}, de l'équation fonctionnelle (1). Mais, comme le montrent les exemples ci-dessous, elle n'appartient pas nécessairement aux solutions définies sur les nombres complexes. Hamel raisonne ainsi. Si $\varphi(x)$ diffère de $C\cdot x$, il existe au moins deux valeurs de la base linéaire\footnote{La base linéaire de tous les nombres (réels) est donnée par un système de nombres (réels) qui permet d'exprimer tous les nombres (réels) linéairement, à coefficients rationnels, tandis qu'entre un nombre fini quelconque de nombres de base il n'existe aucune relation linéaire à coefficients rationnels.}, disons $x_1$ et $x_2$, telles que $\varphi(x_1):x_1\ne\varphi(x_2):x_2$. Alors les deux équations
% unit: noether.p10.eq.011
\[
\alpha x_1+\beta x_2=a;\qquad \alpha\varphi(x_1)+\beta\varphi(x_2)=b
\]
ont une solution en nombres réels $\alpha,\beta$; par conséquent $a$ et $b$ peuvent être approchés arbitrairement par des nombres rationnels $\alpha,\beta$. Comme $\alpha,\beta$ sont rationnels, $\alpha\varphi(x_1)+\beta\varphi(x_2)$ représente alors effectivement $\varphi(\alpha x_1+\beta x_2)$. Ce raisonnement échoue pour les systèmes de valeurs complexes. En effet, on peut donner, aussi dans le domaine complexe, des solutions de (1) qui sont discontinues mais non extrêmement discontinues; par exemple
% unit: noether.p10.eq.012
\[
\varphi(z)=\varphi(x+iy)=\varphi(x)+i(y+\varphi(x)),
\]
où $\varphi(x)$ désigne une solution réelle extrêmement discontinue dans le réel; ou, plus simplement encore, $\varphi(z)=\varphi(x)$. Dans les deux cas la partie réelle de $\varphi(z)$ est extrêmement discontinue, tandis que la partie imaginaire est déterminée par la partie réelle.

% unit: noether.p10.par.016
Ce comportement, et en même temps le comportement de la fonction $f(z)$, devient tout à fait clair si l'on introduit la notion de rang. Posons $z=x+iy$, $\varphi(z)=X+iY$, et appelons $\varphi(z)$ de rang quatre, trois, deux ou un, respectivement, selon qu'il n'existe aucune, une, deux ou trois relations linéairement indépendantes
% unit: noether.p10.eq.013
\[
c_1x+c_2y+c_3X+c_4Y=0,
\]
où les $c$ sont naturellement des nombres réels.

% unit: noether.p10.par.017
1) Que $\varphi(z)$ soit de \emph{rang quatre}. Il existe alors au moins quatre valeurs de la base linéaire, disons $z_1,z_2,z_3,z_4$, telles que le déterminant
% unit: noether.p10.eq.014
\[
\begin{vmatrix}
x_1&x_2&x_3&x_4\\
y_1&y_2&y_3&y_4\\
X_1&X_2&X_3&X_4\\
Y_1&Y_2&Y_3&Y_4
\end{vmatrix}
\]
ne s'annule pas. Les équations
% unit: noether.p10.eq.015
\[
\begin{aligned}
\alpha x_1+\beta x_2+\gamma x_3+\delta x_4&=a,\\
\alpha y_1+\cdots+\delta y_4&=b,\\
\alpha X_1+\cdots+\delta X_4&=A,\\
\alpha Y_1+\cdots+\delta Y_4&=B
\end{aligned}
\]
ont donc une solution en nombres réels $\alpha,\ldots,\delta$, et $a,b,A,B$ peuvent être approchés arbitrairement par des nombres rationnels $\alpha,\beta,\gamma,\delta$. De plus on a
% unit: noether.p10.eq.016
\[
\alpha(X_1+iY_1)+\beta(X_2+iY_2)+\gamma(X_3+iY_3)+\delta(X_4+iY_4)
 =\varphi(\alpha z_1+\beta z_2+\gamma z_3+\delta z_4).
\]
\og{}En général\fg{}, par conséquent, les solutions $\varphi(z)$ sont extrêmement discontinues aussi dans le domaine complexe; les \emph{valeurs de discontinuité} de $\varphi(z)$ au voisinage de $z_0=a+ib$ \emph{forment une variété linéaire bidimensionnelle}.

% unit: noether.p10.par.018
2) Que $\varphi(z)$ soit de \emph{rang trois}. Alors il existe au moins trois systèmes de valeurs de la base linéaire, disons $z_1,z_2,z_3$, tels que le déterminant ci-dessus ait rang trois. Les valeurs $a,b,A,B$ satisfaisant à la relation
% unit: noether.p10.eq.017
\[
c_1a+c_2b+c_3A+c_4B=0,
\]
et celles-là seulement peuvent être approchées arbitrairement; les \emph{valeurs de discontinuité} forment une \emph{variété linéaire unidimensionnelle} déterminée par $z_0=a+ib$. Comme le montrent les exemples indiqués, ce cas peut effectivement se produire.

% unit: noether.p10.par.019
3) Que $\varphi(z)$ soit de \emph{rang deux}. Puisqu'on peut toujours choisir deux nombres complexes $z_1,z_2$ tels que
% unit: noether.p10.eq.018
\[
\begin{vmatrix}x_1&x_2\\ y_1&y_2\end{vmatrix}\ne0,
\]
les relations ont la forme
% unit: noether.p10.eq.019
\[
X=\lambda_1x+\lambda_2y;\qquad Y=\mu_1x+\mu_2y.
\]
L'expression qui en résulte
% unit: noether.p10.eq.020
\[
\varphi(z)=(\lambda_1x+\lambda_2y)+i(\mu_1x+\mu_2y)
\]
est précisément la solution \emph{continue} la plus générale de l'équation fonctionnelle (1) dans le domaine des nombres complexes.

% unit: noether.p10.par.020
4) Que $\varphi(z)$ soit de \emph{rang un}. Ce cas ne peut pas se produire dans le domaine des nombres complexes; dans le domaine des nombres réels il signifie la solution \emph{continue} $\varphi(x)=C\cdot x$.

% unit: noether.p10.par.021
Nous montrons maintenant que les \emph{solutions} $f(z)$ \emph{des équations fonctionnelles (1) à (4), définies pour tous les nombres réels et complexes, ne peuvent avoir que le rang deux ou le rang quatre}. Ici le rang deux ne donne que les deux fonctions \emph{continues} $f(z)=z$ ou $\overline z$, comme le montre la spécialisation $f(1)=1$ et $f(i)=\pm i$, combinée avec 3). Le rang un étant déjà exclu, il reste à exclure le rang trois. Nous supposons donc qu'il existe au moins une relation
% unit: noether.p10.eq.021
\[
c_1x+c_2y+c_3X+c_4Y=0,
\]
de sorte que $f(z)$ ait rang trois, ou peut-être un rang plus bas, et nous montrons que, sous cette hypothèse, c'est toujours le dernier cas qui se produit. Comme le montre de nouveau la spécialisation $f(1)=1$, $f(i)=\pm i$, la relation devient
% unit: noether.p10.eq.022
\[
c_3(X-x)+c_4(Y\mp y)=0,
\]
où $c_3$ et $c_4$ ne s'annulent pas simultanément. Supposons d'abord, par exemple, $c_4=0$. La relation $(X-x)=0$ signifie alors qu'à un $z$ purement imaginaire correspond aussi un $f(z)$ purement imaginaire. Ainsi, pour tout $y$ réel, on a $f(iy)=i\cdot\Yreal$, où $\Yreal$ est de nouveau réel. De là, puisque
% unit: noether.p10.eq.023
\[
f(iy)=\pm i f(y),
\]
on obtient aussi $f(y)=\pm\Yreal$; par conséquent, à tout $z$ réel correspond aussi un $f(z)$ réel, et, à cause de $(X-x)$, toutes les valeurs réelles passent en elles-mêmes. On n'a donc que $f(z)=z$ ou $\overline z$, c'est-à-dire effectivement le rang deux. On conclut de façon tout à fait analogue lorsque $c_3=0$, $c_4\ne0$.

% unit: noether.p10.par.022
Supposons maintenant $c_3$ et $c_4$ différents de zéro, de sorte que la relation puisse s'écrire sous la forme
% unit: noether.p10.eq.024
\[
(Y\mp y)=c\cdot(X-x)
\]
avec $c$ réel, non nul et fini. Cela signifie que pour $y=\pm cx$ on a aussi $Y=c\cdot X$. Ainsi, en tenant compte des équations fonctionnelles, pour tout $x$ réel on a
% unit: noether.p10.eq.025
\[
f(x\cdot(1\pm ic))=\Xreal(1+ic)=f(x)\cdot(1+i f(c)),
\]
où $\Xreal$ est de nouveau réel. Mais ici, d'après le n\textsuperscript{o} 1, le facteur de $f(x)$ ne peut pas s'annuler identiquement, et la division donne
% unit: noether.p10.eq.026
\[
f(x)=\Xreal(\sigma+i\tau)
\]
avec $\sigma$ et $\tau$ réels, indépendants de $x$ et de $\Xreal$. En particulier, à $x=1$ appartient aussi un $X_0$ réel, et la condition $1=X_0(\sigma+i\tau)$ montre que nécessairement $\tau=0$, de sorte que $f(x)=\sigma\cdot\Xreal$. Donc à tout $z$ réel correspond aussi un $f(z)$ réel; de façon équivalente, $y=0$ implique nécessairement $Y=0$. Par là la relation linéaire pour $z$ réel devient $X-x=0$;\footnote{On utilise ici le fait que $c\ne0$, de sorte que ni $c_3$ ni $c_4$ ne s'annule, tandis qu'auparavant seul $c_4\ne0$ avait été utilisé; il fallait donc traiter d'avance le cas où l'un de ces deux s'annule.} chaque valeur réelle correspond à elle-même. On a de nouveau seulement $f(z)=z$ ou $\overline z$, donc effectivement le rang deux. Le rang trois est ainsi exclu dans tous les cas. Mais, puisque --- comme le montrent les recherches des premiers numéros --- des solutions discontinues existent effectivement et doivent donc nécessairement avoir le rang quatre, on obtient d'après 1) le théorème: \emph{\og{}Toutes les solutions} $f(z)$ \emph{des équations fonctionnelles (1) à (4) --- à l'exception des solutions continues} $f(z)=z$ \emph{ou} $\overline z$ --- \emph{sont extrêmement discontinues dans le réel et dans le complexe}.\fg{}\footnote{Lebesgue montre loc. cit. que la partie réelle et la partie imaginaire de $f(z)$ sont toutes deux des fonctions \og{}non mesurables\fg{} de $x$ et de $y$; comme le montre l'exemple du début du n\textsuperscript{o} 4, $\varphi(z)=\varphi(x)+i(y+\varphi(x))$, on ne peut pas encore en conclure l'extrême discontinuité dans le domaine complexe.}

% unit: noether.p10.close.001
\medskip
\noindent Erlangen, 30 octobre 1915.
